\chapter{SYSTEM STUDY AND FEASIBILITY ANALYSIS}
\label{ch:system_study}

\section{Existing System Study}

\subsection{Workflow of Existing Heuristic Speech Distillation}
In conventional knowledge distillation pipelines designed for speech recognition decoders, transferring linguistic priors from large autoregressive language models to compact acoustic front-ends relies on discrete string-matching heuristics. In these frameworks, the teacher model processes the ground-truth text sequence to generate soft probability distributions over its high-capacity subword vocabulary. Concurrently, the student acoustic model maps speech spectrogram frames into a significantly smaller, domain-specific vocabulary.

Because the vocabulary indices between teacher and student do not correspond to identical lexical units ($\mathcal{V}_{\text{teacher}} \ne \mathcal{V}_{\text{student}}$), classical Kullback-Leibler (KL) divergence cannot be computed directly across raw output logits. To circumvent this structural barrier, existing frameworks implement a sequential edit-distance alignment stage:
\begin{enumerate}
    \item \textbf{Token-to-Token String Mapping:} The system evaluates the normalized Levenshtein string distance $d(u_i, v_j)$ between subword strings generated by the teacher and candidate subwords of the student.
    \item \textbf{Heuristic Discard Criterion:} A hard distance threshold is enforced (typically $d \le 0.2$). If a student subword matches a teacher subword within this threshold, a one-to-one pseudo-ground-truth mapping is established.
    \item \textbf{Logit Rescaling & Partial Loss Computation:} For matched tokens, cross-entropy or KL divergence loss is calculated, while candidate tokens exceeding the distance threshold are completely omitted from the distillation gradient update.
\end{enumerate}

\subsection{Deficiencies of the Existing Approach}
While heuristic string distillation provides a preliminary bridge between language models and speech decoders, comprehensive architectural analysis reveals four critical failure modes:
\begin{enumerate}
    \item \textbf{Token Discard and Information Loss:} Natural language subword tokenizers decompose words into arbitrary morphological fragments based on training frequency. For example, a teacher model with a 128k vocabulary may tokenize the word ``observing'' as a single discrete token, whereas a compact student with a 1k vocabulary splits it into ``obs'' and ``erving''. Under rigid edit-distance thresholds, multi-token fragments exhibit high string distance ($d > 0.2$) and are discarded entirely from supervision. Consequently, the student receives zero gradient guidance on complex or polysyllabic vocabulary.
    \item \textbf{Narrow Candidate Truncation Overhead:} Calculating string edit distances dynamically across millions of sequence permutations in real-time training batches incurs severe computational latency. To maintain tractable batch times, existing systems truncate supervision to the top-512 candidate tokens, ignoring the tail of the semantic probability distribution where subtle syntactic priors reside.
    \item \textbf{Absence of Paralinguistic & Affective Modeling:} Existing speech distillation architectures operate solely on linguistic transcript matching. Acoustic modulations, emotional prosody, and affective speaker states are ignored, preventing the framework from being utilized in emotion-aware conversational applications.
    \item \textbf{Decoupled Multi-Stage Execution:} Conventional setups maintain frozen acoustic feature extractors (such as 95M-parameter self-supervised front-ends) coupled to external autoregressive decoders, resulting in total active parameter counts exceeding 130M weights and impeding low-latency edge deployment.
\end{enumerate}

\section{Proposed System Concept}
The proposed framework supersedes heuristic string alignment by establishing a continuous, geometric distillation bridge governed by entropic Optimal Transport. Rather than evaluating discrete token identity matches, our approach projects both teacher and student representations into a continuous, shared metric embedding space. Probability mass is transported smoothly from teacher distributions to student decoders based on continuous semantic distances, guaranteeing that no subwords are discarded, regardless of vocabulary cardinality disparities.

Furthermore, the proposed system introduces a unified multi-task architecture: a single, ultra-compact 12.55M-parameter Conformer student network simultaneously drives non-autoregressive phonetic transcription under Connectionist Temporal Classification (CTC) and 4-class categorical Speech Emotion Recognition (SER) via temporal mean pooling. The causal foundation teacher remains strictly static throughout training, incurring zero computational or memory overhead during test-time inference.

\section{Feasibility Analysis}

\subsection{Technical Feasibility}
Technical feasibility evaluates whether the proposed multi-task distillation framework can be trained within accessible academic compute constraints and deployed on resource-constrained embedded edge hardware.
\begin{enumerate}
    \item \textbf{Inference Memory Budget:} The complete student network comprises only 12.55 million active parameters. In standard 32-bit floating-point format, the entire network requires only 50.2 MB of resident memory. This fits comfortably within the strict RAM constraints of embedded edge processors, mobile microcontrollers, and single-board computers (which typically provide 1 GB to 4 GB of host memory).
    \item \textbf{Real-Time Latency Feasibility:} Non-autoregressive CTC decoding eliminates the sequential autoregressive loop, allowing all acoustic frames to be decoded in parallel. Operating at an edge Real-Time Factor of $\text{RTF}_{\text{CPU}} = 0.08$ on commodity multi-core processors, the model processes a 10-second audio utterance in just 0.8 seconds—operating 12.5$\times$ faster than real-time incoming speech.
    \item \textbf{Training Memory Management:} While exact Optimal Transport solved via linear programming requires cubic complexity $\mathcal{O}(N^3)$, incorporating entropic regularization solved via the Sinkhorn-Knopp matrix scaling algorithm reduces computational complexity to matrix-vector multiplications scaling as $\mathcal{O}(M \cdot N)$. This allows training batches to execute efficiently within consumer-grade accelerator memory without memory exhaustion.
\end{enumerate}

\subsection{Operational Feasibility}
Operational feasibility assesses the usability, integration potential, and practical deployment viability of the developed speech interface:
\begin{itemize}
    \item \textbf{Universal Hands-Free Edge Interfaces:} Voice-controlled Internet-of-Things (IoT) appliances, automotive dashboards, and medical dictation units require speech interfaces that operate locally without transmitting raw audio to cloud servers, preserving privacy and ensuring offline functionality.
    \item \textbf{Empathetic Spoken Dialog Systems:} Incorporating real-time emotion recognition alongside speech transcription enables conversational systems to detect distress, frustration, or satisfaction in customer service interactions, mental health monitoring, and interactive learning environments.
    \item \textbf{Zero-Configuration Deployment:} Because the heavy foundation teacher model is completely decoupled from the student network during distillation, end-users require only the lightweight student weights, completely eliminating the need for complex multi-model server setups.
\end{itemize}

\subsection{Economic and Resource Feasibility}
From an economic standpoint, deploying commercial LLM-based speech services in enterprise environments incurs exorbitant cloud API subscription costs, high electricity consumption, and continuous bandwidth overhead. The proposed solution delivers substantial economic advantages:
\begin{itemize}
    \item \textbf{Elimination of Cloud Server Dependencies:} By transferring the linguistic intelligence of a billion-parameter foundation model into a 12.55M edge model, high-accuracy speech processing is executed directly on local hardware, reducing ongoing cloud inference expenses to zero.
    \item \textbf{Academic Resource Accessibility:} The entire training pipeline is engineered to converge on consumer-grade accelerator hardware, demonstrating that state-of-the-art multimodal foundation distillation does not require industrial-scale supercomputing clusters.
\end{itemize}

\section{Dataset Feasibility and Corpus Profiles}

\subsection{LibriSpeech 960-Hour Audiobook Corpus}
For acoustic backbone bootstrapping and linguistic pre-training, the framework utilizes the widely recognized LibriSpeech corpus \cite{panayotov2015librispeech}. Derived from read public-domain audiobooks, LibriSpeech provides clean, standardized English speech sampled at 16 kHz mono format. 

Key attributes of the LibriSpeech benchmark include:
\begin{itemize}
    \item \textbf{Acoustic Quality:} Clean studio-grade recordings characterized by high signal-to-noise ratios, clear phonetic articulation, and consistent recording environments.
    \item \textbf{Speaker Diversity:} Encompasses approximately 2,500 distinct speakers with a balanced distribution of male and female voices across diverse English accents.
    \item \textbf{Subword Segmentation:} Transcripts are tokenized using a unigram SentencePiece Byte-Pair Encoding (BPE) model with a vocabulary size of $|\mathcal{V}_{\text{student}}| = 1{,}024$ tokens, reserving dedicated indices for special control tokens.
\end{itemize}

\subsection{IEMOCAP Multimodal Emotional Database}
To evaluate multimodal speech emotion recognition alongside acoustic transcription, we employ the Interactive Emotional Dyadic Motion Capture (IEMOCAP) database \cite{busso2008iemocap}. IEMOCAP contains approximately 12 hours of synchronized audiovisual speech data collected from 10 professional actors (5 male, 5 female) organized across 5 distinct dyadic recording sessions.

Table \ref{tab:iemocap_stats} details the categorical distribution of emotional utterances within the IEMOCAP corpus.

\begin{table}[htbp]
\centering
\caption{Categorical and Speaker Distribution in the IEMOCAP Corpus}
\label{tab:iemocap_stats}
\begin{tabular}{lccc}
\toprule
\textbf{Emotion Category} & \textbf{Utterance Count} & \textbf{Percentage (\%)} & \textbf{Evaluation Split} \\
\midrule
Neutral & 1,708 & 16.3\% & 4-Class Core Benchmark \\
Happy / Excited & 1,636 & 15.6\% & 4-Class Core Benchmark \\
Sad & 1,084 & 10.3\% & 4-Class Core Benchmark \\
Angry & 1,103 & 10.5\% & 4-Class Core Benchmark \\
Frustrated / Disputed & 4,863 & 46.3\% & Non-Consensus / Ambiguous \\
\midrule
\textbf{Total Dataset} & \textbf{10,494} & \textbf{100.0\%} & \textbf{5 Dyadic Sessions (10 Actors)} \\
\bottomrule
\end{tabular}
\end{table}

\subsection{The Non-Consensus Dilemma: 2,507 Ambiguous Samples}
A critical phenomenon in affective corpus curation is human evaluator disagreement. In the IEMOCAP database, each dialogue turn is annotated by three independent human evaluators. While 5,531 utterances attain majority consensus (at least 2 out of 3 evaluators agreeing on a single category), exactly 2,507 utterances reside in a non-consensus state, designated as the ambiguous class (\texttt{xxx}). 

In conventional academic practice, these 2,507 samples are discarded as unusable noise. In this project, we retain these samples to evaluate whether the distilled multimodal representation can serve as an objective mathematical tie-breaker, evaluating Human-AI consensus across real-world perceptual ambiguity.

\section{SWOT Analysis of the Framework}
To provide a structured evaluation of the proposed distillation methodology, a comprehensive SWOT (Strengths, Weaknesses, Opportunities, and Threats) analysis is conducted.

\subsection{Strengths}
\begin{itemize}
    \item \textbf{Continuous Geometric Mapping:} Replaces brittle string-match thresholds with continuous cosine embedding distances, ensuring 100\% token retention across disparate subword vocabularies.
    \item \textbf{Mitigation of CTC Blank Collapse:} Subsampling acoustic frames by $4\times$ to 25 fps, combined with Sinkhorn semantic grounding, systematically contracts null token dominance from 92.4\% down to 68.2\%.
    \item \textbf{Zero Inference Overhead:} The foundation teacher model is completely detached during inference, allowing the 12.55M student to operate on low-power edge CPUs at an RTF of 0.08.
\end{itemize}

\subsection{Weaknesses}
\begin{itemize}
    \item \textbf{Iterative Scaling Complexity during Training:} Although faster than linear programming, Sinkhorn-Knopp matrix scaling requires multiple forward-backward vector-matrix iterations per batch, marginally extending training epoch duration.
    \item \textbf{Reliance on Pretrained Embedding Quality:} The transportation plan is fundamentally governed by the semantic structure of the teacher token embeddings; poorly clustered embedding manifolds could introduce noisy transport gradients.
\end{itemize}

\subsection{Opportunities}
\begin{itemize}
    \item \textbf{On-Device Edge Deployment:} Enables offline speech recognition and empathetic conversational assistance on smart home appliances, wearable robotics, and hearing aids.
    \item \textbf{Multilingual & Cross-Lingual Distillation:} The geometric transport formulation can be extended to transfer linguistic priors from multilingual foundation models (such as multilingual LLMs) to low-resource regional dialects.
\end{itemize}

\subsection{Threats}
\begin{itemize}
    \item \textbf{Numerical Instability under Small Entropic Bandwidths:} Setting the regularization parameter too small ($\gamma < 0.01$) induces floating-point underflow in exponential kernel computations.
    \item \textbf{Task Gradient Interference:} Over-weighting the affective classification objective ($\lambda_{\text{emo}} > 0.5$) can degrade phoneme boundary detection, requiring calibrated multi-task loss scheduling.
\end{itemize}
